\documentclass[final]{beamer}

% ====================
% Packages
% ====================

\usepackage[T1]{fontenc}
\usepackage[utf8]{luainputenc}
\usepackage{lmodern}
\usepackage[size=custom, width=122,height=91, scale=1.2]{beamerposter}
\usetheme{gemini}
\usecolortheme{msu}
\usepackage{graphicx}
\usepackage{booktabs}
\usepackage{tikz}
\usepackage{pgfplots}
\pgfplotsset{compat=1.14}
\usepackage{anyfontsize}
\usepackage{tabularx}

% Override exampleblock background to light yellow
\definecolor{lightyellow}{RGB}{255, 252, 215}
\setbeamercolor{block example title}{fg=umichblue,bg=lightyellow}
\setbeamercolor{block example separator}{bg=black}
\setbeamercolor{block example body}{fg=black,bg=lightyellow}

% ====================
% Lengths
% ====================

\newlength{\sepwidth}
\newlength{\colwidth}
\setlength{\sepwidth}{0.025\paperwidth}
\setlength{\colwidth}{0.3\paperwidth}

\newcommand{\separatorcolumn}{\begin{column}{\sepwidth}\end{column}}

% ====================
% Title
% ====================

\title{From Sentiment to Signals: FinMini LLM for Financial Sentiment Analysis and Text-Driven Stock Movement Prediction}

\author{JiaYi Zhang \and JiaYuan Jiang \and Yue Yu \and YiXiang Fei}

\institute[UMich]{Electrical Engineering and Computer Science, University of Michigan \quad $\cdot$ \quad EECS 545: Machine Learning}

% ====================
% Footer
% ====================

\footercontent{
  EECS 545 --- Machine Learning \hfill
  University of Michigan, April 2026 \hfill
  \href{mailtojiayuanj@umich.edu}{\texttt{github.com/jiayuan-jiang/EECS545\_project}}}

% ====================
% Logo
% ====================

\logoright{\includegraphics[height=5cm]{logos/blockM.png}}

% ====================
% Body
% ====================

\begin{document}

\begin{frame}[t]
\begin{columns}[t]
\separatorcolumn

% ==================== COLUMN 1 ====================
\begin{column}{\colwidth}

  \begin{block}{Motivation}

    Large Language Models (LLMs) have demonstrated impressive general-purpose
    capabilities, yet they consistently underperform in specialized financial
    domains. Financial texts carry \textbf{domain-specific terminology},
    \textbf{implicit sentiment cues}, and \textbf{complex causal reasoning}
    absent from general corpora --- for example, a phrase like
    ``earnings in line with expectations'' may signal a mildly negative
    outcome to a financial analyst, while a general LLM interprets it as
    neutral. We target two core tasks --- \textbf{financial sentiment
    analysis} and \textbf{text-driven stock movement prediction} --- by
    fine-tuning a compact 8B model with QLoRA on curated financial data,
    and investigate whether CoT supervision further improves financial
    reasoning beyond label-only SFT.

  \end{block}

  \begin{block}{Problem Statement}

    Given a piece of financial text --- such as a news headline, analyst
    commentary, or social media post --- we address two binary
    classification tasks:

    \begin{itemize}
      \item \textbf{Financial Sentiment Analysis:} determine whether the
        text conveys a \emph{positive} or \emph{negative} market-relevant
        signal.
      \item \textbf{Stock Movement Prediction:} predict whether the
        referenced stock will \emph{rise} or \emph{fall} based on textual
        and quantitative signals.
    \end{itemize}

    Both tasks are unified under a single positive/negative framework.
    We investigate a two-stage fine-tuning strategy: Stage~1 adapts the
    base model via label-only SFT; Stage~2 enhances reasoning via
    Chain-of-Thought SFT with structured \texttt{<think>} traces.
    The central hypothesis is that CoT supervision yields more
    interpretable and accurate predictions than label-only fine-tuning,
    particularly on out-of-distribution benchmarks.

  \end{block}

  \begin{block}{Related Work}

    \textbf{FinBERT} (Araci, 2019) established domain-specific pre-training
    as the baseline for financial NLP, showing BERT fine-tuned on financial
    text outperforms general-purpose models on sentiment tasks.

    \textbf{BloombergGPT} (Wu et al., 2023) demonstrated that large-scale
    financial pre-training yields strong performance but requires massive
    proprietary data, motivating parameter-efficient alternatives.

    \textbf{PIXIU} (Xie et al., 2023) introduced a comprehensive
    financial instruction-tuning benchmark (FinBen) and showed that
    supervised fine-tuning on diverse financial tasks achieves competitive
    results on the Open FinLLM Leaderboard --- the evaluation framework
    we adopt.

    \textbf{Long CoT SFT} (Ou \& Yin, 2025) demonstrated that
    chain-of-thought supervised fine-tuning is a critical prerequisite
    for instilling robust reasoning in lightweight LLMs, outperforming
    reinforcement learning as a standalone strategy --- directly
    motivating our Stage~2 CoT-SFT design.

  \end{block}

  \begin{block}{References}

    \scriptsize
    \setlength{\parskip}{0.3ex}
    \noindent
    [1] Ou \& Yin. \textit{Empowering Lightweight MLLMs via Long CoT SFT}. arXiv:2509.03321, 2025.

    \noindent
    [2] Araci. \textit{FinBERT: Financial Sentiment Analysis with Pre-trained LMs}. arXiv:1908.10063, 2019.

    \noindent
    [3] Xie et al. \textit{PIXIU: A Large Language Model and Benchmark for Finance}. NeurIPS, 2023.

    \noindent
    [4] Wu et al. \textit{BloombergGPT: A Large Language Model for Finance}. arXiv:2303.17564, 2023.

    \noindent
    [5] Huang et al. \textit{Open-FinLLMs: Open Multimodal LLMs for Financial Applications}. arXiv:2402.04616, 2024.

  \end{block}

\end{column}

\separatorcolumn

% ==================== COLUMN 2 ====================
\begin{column}{\colwidth}

  \begin{block}{Approach}

    \heading{Two-Stage Fine-Tuning Pipeline}

    \begin{figure}
      \centering
      \includegraphics[width=0.7\textwidth]{pipeline_diagram.png}
    \end{figure}

    \heading{Base Model and Fine-Tuning Method}

    \begin{figure}
      \centering
      \includegraphics[width=0.7\textwidth]{qlora_diagram.png}
    \end{figure}

  \end{block}

  \begin{exampleblock}{Data}

  Fine-tuning quality is bottlenecked by data quality and diversity.
    We select three datasets that balance \textbf{annotation reliability}
    (FPB: 100\% inter-annotator agreement), \textbf{scale} (DS2: 100K+
    multi-source samples), and \textbf{task diversity} (FLARE-SM-ACL:
    market-movement reasoning). Each sample is further augmented via
    \textbf{DeepSeek API} into CoT-SFT format, replacing bare labels
    with structured \texttt{<think>}\ldots\texttt{</think>} reasoning
    traces, producing interpretable, reasoning-rich training signal
    aligned with DeepSeek-R1's native format (Ou \& Yin, 2025).

    \heading{Training Corpus \normalfont\small(63,395 train / 7,979 test)}

    \begin{table}
      \centering\small
      \begin{tabularx}{\linewidth}{@{}lrX@{}}
        \toprule
        \textbf{Dataset} & \textbf{Size} & \textbf{Selection Rationale} \\
        \midrule
        FPB          &   2,264 & 100\% inter-annotator agreement; gold-standard label quality \\[0.3ex]
        DS2          & 100,037 & Multi-source (FinGPT, Twitter); provides training scale \\[0.3ex]
        FLARE-SM-ACL &  27,056 & Stock movement reasoning; extends task diversity \\
        \bottomrule
      \end{tabularx}
    \end{table}

    \heading{Evaluation Benchmarks}

    \begin{table}
      \centering\small
      \begin{tabularx}{\linewidth}{@{}lrX@{}}
        \toprule
        \textbf{Benchmark} & \textbf{Size} & \textbf{Description} \\
        \midrule
        en-FPB     &   393 & FPB (FLARE ver.); in-distribution sentiment reference \\[0.3ex]
        FiQA-SA    &   217 & Financial Q\&A opinions; unseen cross-domain sentiment \\[0.3ex]
        SM-BigData & 1,472 & BigData22 stock movement; unseen social media signals \\[0.3ex]
        SM-CIKM    & 1,143 & CIKM18 stock movement; unseen market prediction \\
        \bottomrule
      \end{tabularx}
    \end{table}


    % Four benchmarks from the Open FinLLM Leaderboard were used:
    % \textbf{en-FPB} as an in-distribution reference, and
    % \textbf{FiQA-SA}, \textbf{SM-BigData}, \textbf{SM-CIKM} as
    % entirely unseen datasets to assess generalization across
    % sentiment analysis and stock movement tasks.

    % \heading{CoT Augmentation via DeepSeek API}

    % Each training sample was augmented by querying \textbf{DeepSeek-chat}
    % (temp$=0.7$, max\_tokens$=300$) to generate a structured reasoning
    % trace. The model output was reformatted from a bare label into a
    % \texttt{<think>}\,\ldots\,\texttt{</think>} reasoning chain followed
    % by the final sentiment label, yielding a \textbf{CoT-SFT} dataset.

    % \begin{itemize}
    %   \item \textbf{Interpretability:} predictions with explicit
    %     financial reasoning, not just a label
    %   \item \textbf{Reasoning capability:} long CoT SFT is a proven
    %     prerequisite for instilling reasoning in lightweight LLMs
    %     (Ou \& Yin, 2025) 
    %   \item \textbf{Format alignment:} matches DeepSeek-R1's native
    %     \texttt{<think>} format, enabling effective adapter fine-tuning
    % \end{itemize}

    All datasets were unified into a binary positive/negative format;
    neutral labels were removed. en-FPB shares the same source as the
    training FPB, while the remaining three benchmarks are entirely
    unseen, testing out-of-distribution generalization.

  \end{exampleblock}

\end{column}

\separatorcolumn

% ==================== COLUMN 3 ====================
\begin{column}{\colwidth}

  \begin{block}{Results and Evaluation}

    \heading{Training Convergence}

    Two-stage training on DeepSeek-R1-Distill-Llama-8B (LoRA $r=16$):
    Stage~1 QLoRA SFT loss drops from ${\sim}4.0$ to ${\sim}1.8$;
    Stage~2 CoT-SFT converges to ${\sim}0.5$, reflecting the richer
    chain-of-thought supervision signal.

    \begin{figure}
      \centering
      \includegraphics[width=0.85\textwidth]{figures/fig8_qlora_vs_cotsft_loss.png}
    \end{figure}

    \heading{Accuracy on Open FinLLM Benchmarks (\%)}

    We evaluate on four benchmarks from the Open FinLLM Leaderboard:
    FPB (news sentiment), FiQASA (Q\&A sentiment), SM-BigData and SM-CIKM
    (stock movement prediction).
    FinMini achieves the best accuracy on all four tasks.
    Gains are largest on sentiment benchmarks (+9.6\% on FPB, +12.5\% on FiQASA
    over the strongest baseline), while stock movement tasks also improve,
    confirming that domain fine-tuning generalizes across financial NLP.

    % \begin{table}
    %   \centering
    %   \begin{tabular}{lcccc}
    %     \toprule
    %     \textbf{Model} & \textbf{FPB} & \textbf{FiQASA} & \textbf{SM-BD} & \textbf{SM-CIKM} \\
    %     \midrule
    %     Mistral-7B              & 85.5 & 72.8 & 49.1 & 47.3 \\
    %     Llama-3-8B              & 69.5 & 82.0 & 43.1 & 39.3 \\
    %     DeepSeek-8B (base)      & 83.5 & 57.6 & 55.0 & 49.3 \\
    %     \textbf{FinMini (Ours)} & \textbf{93.1} & \textbf{94.5} & \textbf{57.1} & \textbf{52.2} \\
    %     \bottomrule
    %   \end{tabular}
    % \end{table}
    \begin{figure}
      \centering
      \includegraphics[width=\textwidth]{figures/fig7_bar.png}
    \end{figure}

  \end{block}

  \begin{alertblock}{Conclusion}

    \begin{itemize}
      \item QLoRA SFT on domain-specific data consistently outperforms
        general-purpose baselines across all four benchmarks, with
        \textbf{+9.6\%} on FPB and \textbf{+12.5\%} on FiQASA over the
        strongest baseline.
      \item CoT-SFT further boosts sentiment accuracy to
        \textbf{98.5\%} (FPB) and \textbf{97.7\%} (FiQASA), confirming
        that structured reasoning supervision significantly enhances
        financial NLP performance.
      \item Both stages improve stock movement prediction, demonstrating
        that domain fine-tuning generalizes beyond pure sentiment tasks.
    \end{itemize}

  \end{alertblock}


\end{column}

\separatorcolumn
\end{columns}
\end{frame}

\end{document}
